\section{指数族把归一化、矩与似然连接起来}
\label{sec:13-Machine-Learning/06-Probabilistic-Models/06}

你在一个建模库里接连实现了几个估计器。Bernoulli 的最大似然是数出 \(1\) 的比例；Poisson 的是样本均值；已知方差的 Gaussian 还是样本均值；GMM 的 M 步是责任加权的均值与二阶矩；HMM 的发射参数更新也是同一副样子。三份推导各写一遍，得到的代码几乎可以互相复制粘贴——都是先把数据压成几个求和，再令模型的某个期望等于它。

反复出现的结构值得怀疑。这些分布的公式看上去毫无关系，为什么估计方程会长成同一个模样？答案不在任何一个分布里，而在它们共同所属的那个族。

\subsection{一个族的形状由充分统计量和一个归一化函数定下}

指数族把上述共同骨架写成

\[
p(x\given\eta)=h(x)\exp\bigl(\eta^{\trans}T(x)-A(\eta)\bigr),
\]

其中 \(T(x)\) 是充分统计量，\(\eta\) 是自然参数，\(h(x)\) 是与参数无关的基测度，而

\[
A(\eta)=\log\int h(x)\exp\bigl(\eta^{\trans}T(x)\bigr)\,dx
\]

是 log-partition function（离散情形把积分换成求和）。\(A\) 不是为了让式子好看才补上的常数，它要保证密度积分为 \(1\)，并且——这是本节的主线——它同时装着这个族的全部矩信息。

允许的自然参数只能取在 \(\mathcal H=\set{\eta:A(\eta)<\infty}\) 内，并非任意向量都合法。均值与方差都未知的一维 Gaussian 可以写成

\[
T(x)=\begin{bmatrix}x\\x^2\end{bmatrix},
\qquad
\eta=\begin{bmatrix}\mu/\sigma^2\\-1/(2\sigma^2)\end{bmatrix},
\]

第二个分量必须严格为负，否则 \(\exp(\eta_2x^2)\) 在两端发散，积分不存在。自然参数空间的边界因此对应着方差退化、概率彻底集中、或者最大似然解跑掉这类退化情形，后面会看到具体例子。

参数化本身也需要一点小心。若 \(T(x)\) 的各分量之间存在线性冗余，不同的 \(\eta\) 会给出同一个分布，这叫非最小表示。换成最小表示之后，参数与分布一一对应，下面关于严格凸性的说法才成立。

\subsection{Log-partition 求一次导是均值，求两次导是协方差}

在可以交换求导与积分的条件下，直接对 \(A\) 求梯度：

\[
\nabla_\eta A(\eta)
=\frac{\int T(x)h(x)\exp(\eta^{\trans}T(x))\,dx}
{\int h(x)\exp(\eta^{\trans}T(x))\,dx}
=\E_\eta[T(X)].
\]

分母正是配分函数，分子把 \(T(x)\) 带进同一个积分，商就是期望。再求一次导，

\[
\nabla_\eta^2A(\eta)=\Cov_\eta\bigl(T(X)\bigr).
\]

协方差矩阵半正定，所以 \(A\) 是凸函数；最小表示的正则指数族中它严格凸。于是归一化、期望与曲率被同一个函数管起来：\(\nabla A\) 把自然参数映到均值参数，\(\nabla^2A\) 描述统计量对参数变动的敏感程度，也正是 Fisher 信息，可对照\hyperref[sec:11-Mathematical-Statistics/03-Sufficient-Statistics-and-Information/02]{``Score 与 Fisher 信息衡量局部可辨认性''}。

以 Bernoulli 为例把这件事画出来。写成指数族形式 \(p(x\given\eta)=\exp(\eta x-\log(1+e^\eta))\)，\(x\in\set{0,1}\)，自然参数就是 log-odds，\(A(\eta)=\log(1+e^\eta)\)，于是 \(\nabla A(\eta)=\sigma(\eta)\)。

\begin{center}
\begin{tikzpicture}[>=stealth]

% ---------- 左：A(eta) ----------
\begin{scope}
\draw[gray!60,->] (-3.4,0) -- (3.6,0);
\draw[gray!60,->] (0,-0.25) -- (0,3.7);
\node[font=\scriptsize,anchor=north] at (3.5,-0.05) {\(\eta\)};

\draw[black!30,dashed] (0,0) -- (3.4,3.4);
\draw[violet,very thick] plot[domain=-3.3:3.3,samples=90]
  (\x,{ln(1+exp(\x))});
\draw[blue!70,thick] (-0.6,0.146) -- (3.0,2.775);
\fill[blue!80] (1,1.313) circle (2pt);

\node[font=\scriptsize,violet,anchor=west] at (0.15,3.35) {\(A(\eta)=\log(1+e^{\eta})\)};
\node[font=\scriptsize,blue!70,anchor=west] at (1.5,0.95) {斜率 \(=\E_\eta[T(X)]\)};
\node[font=\scriptsize,anchor=north] at (0.2,-0.15) {自然参数};
\end{scope}

% ---------- 右：均值映射 ----------
\begin{scope}[shift={(7.6,0)}]
\draw[gray!60,->] (-3.4,0) -- (3.6,0);
\draw[gray!60,->] (0,-0.25) -- (0,3.2);
\node[font=\scriptsize,anchor=north] at (3.5,-0.05) {\(\eta\)};

\draw[black!35,dashed] (-3.3,2.5) -- (3.3,2.5);
\node[font=\scriptsize,anchor=east] at (-0.08,2.5) {\(1\)};
\node[font=\scriptsize,anchor=east] at (-0.08,0.18) {\(0\)};

\draw[violet,very thick] plot[domain=-3.3:3.3,samples=90]
  (\x,{2.5/(1+exp(-\x))});
\fill[blue!80] (1,1.828) circle (2pt);
\draw[gray!55,dashed] (1,0) -- (1,1.828);

\node[font=\scriptsize,anchor=west] at (-3.3,2.85) {\(\nabla A(\eta)=\sigma(\eta)\)};
\node[font=\scriptsize,anchor=east] at (3.3,1.6) {够不到边界};
\end{scope}

\end{tikzpicture}
\end{center}

\emph{左边曲线的斜率就是右边的高度；边界那条虚线只能逼近，不能落上去。}

左图那条切线的斜率读出的是分布的均值，这正是 \(\nabla A=\E[T]\) 的几何含义；曲线越弯，同样的参数扰动带来的均值变化越大，对应 \(\nabla^2A=\Cov(T)\)。右图是同一件事的另一种画法：均值映射把整条实轴压进开区间 \((0,1)\)，两端的虚线够不到。这张图稍后会解释一个常见的训练故障。

同样的代数在\hyperref[sec:13-Machine-Learning/08-Sampling-and-Inference/06]{``配分函数把能量变成全局概率''}里以能量模型的面目出现。区别在于指数族把能量限制成统计量的线性组合，配分函数的梯度因而退化成一个明确的矩映射，而不是一般能量模型里那个难算的期望。最大熵原理从另一头给出同一族分布：在若干矩约束下熵最大的分布必是指数族形式，见\hyperref[sec:07-Information-Theory/01-Information-and-Entropy/04]{``最大熵原理：从约束到分布''}。

\subsection{独立样本只通过充分统计量进入似然}

对独立样本 \(x_{1:n}\)，联合密度是

\[
p(x_{1:n}\given\eta)
=\Bigl(\prod_{i=1}^nh(x_i)\Bigr)
\exp\Bigl(\eta^{\trans}\sum_{i=1}^nT(x_i)-nA(\eta)\Bigr).
\]

与参数有关的部分只碰到 \(S_n=\sum_iT(x_i)\)。按 Fisher--Neyman 分解，\(S_n\) 是 \(\eta\) 的充分统计量：不论样本量多大，参数推断需要的东西被压进一个固定维数的向量。要注意``充分''的确切含义——给定 \(S_n\) 后原始数据对参数不再提供额外的似然信息，但它对别的问题仍可能重要。样本顺序、离群结构、模型是否失配，这些都在压缩中消失了，所以充分统计量不能拿来做模型诊断，边界见\hyperref[sec:11-Mathematical-Statistics/03-Sufficient-Statistics-and-Information/01]{``充分统计量保留了哪些参数信息''}。

现在把开头那个疑问收掉。对数似然的梯度为

\[
\nabla_\eta\ell(\eta)=\sum_{i=1}^nT(x_i)-n\nabla A(\eta),
\]

令其为零，内点上的最大似然解满足

\[
\frac1n\sum_{i=1}^nT(x_i)=\E_{\widehat\eta}[T(X)].
\]

经验矩等于模型矩——矩匹配不是某种巧合或者启发式，它就是指数族似然的一阶条件。Bernoulli 的 \(T(x)=x\) 于是给出样本比例，Gaussian 的 \(T(x)=(x,x^2)^{\trans}\) 给出样本均值与二阶矩，GMM 的 M 步在责任加权之后仍是这套方程。Hessian 是

\[
\nabla_\eta^2\ell(\eta)=-n\Cov_\eta\bigl(T(X)\bigr),
\]

半负定，所以最小表示下的对数似然是凹的，局部极大即全局极大（凹性的判据见\hyperref[sec:09-Convex-Optimization/03-Convex-Functions/01]{``弦在图像上方意味着什么''}）。这也解释了 EM 的 M 步为什么总能闭式：完全数据的对数似然若属于指数族，M 步不过是在加权的充分统计量上再做一次矩匹配。

\subsection{当经验矩落到边界上，MLE 就不存在}

回到右图。若样本里既有 \(0\) 又有 \(1\)，经验均值落在开区间内部，矩匹配给出一个有限的 \(\widehat\eta\)。若全部观测都是 \(1\)，经验均值等于 \(1\)，恰好落在均值参数空间的边界——那条虚线上没有任何 \(\eta\) 的像。似然沿 \(\eta\to+\infty\) 单调上升却永不取到上确界，有限的最大似然解不存在。

Logistic 回归里完全可分的数据是同一件事的高维版本：存在一个方向，沿它放大系数能让似然持续增加，优化器于是一路把范数推大。日志上看到系数发散、迭代不收敛，第一反应不该是换优化器或调步长，而是检查经验矩是不是贴上了边界。加先验、加 \(\ell_2\) 正则、或者换一个不那么饱和的模型都能把解拉回有限区域，代价是估计的对象已经不再是原来那个。

\subsection{指数族不是``常见分布''的同义词}

这套代数好用，但它的适用范围比印象中窄。有限维指数族要求样本量增加时充分统计量的维数保持固定，Uniform 这种支撑依赖参数的分布不满足正则条件，重尾分布往往写不成这个形式。混合模型更值得单独一提：每个成分都是指数族，隐变量边缘掉之后的观测似然却含着对数里的求和，凹性随之丢失——前面几节反复遇到的局部最优问题，不会因为成分属于指数族而消失。

即便能写成指数形式，也要看自然参数空间是否开、是否光滑，否则上面关于导数与凹性的推理逐条失效。把这些条件核对清楚之后，指数族给出的回报是相当大的：一个统一的估计方程、一个凹的目标、一个把参数几何与信息几何连起来的入口（见\hyperref[sec:12-Real-Analysis-Support/09-Information-Geometry/01]{``统计流形把参数族变成几何对象''}），以及广义线性模型的整套框架（见\hyperref[sec:13-Machine-Learning/02-Linear-Models/05]{``GLM 用链接函数连接均值与线性预测子''}）。还有一项回报尚未兑现：如果似然属于指数族，先验该怎么选，才能让 Bayes 更新之后的后验仍然待在同一族里。下一节处理这个问题。
